%Version 3.1 December 2024
% See section 11 of the User Manual for version history
%
%%%%%%%%%%%%%%%%%%%%%%%%%%%%%%%%%%%%%%%%%%%%%%%%%%%%%%%%%%%%%%%%%%%%%%
%%                                                                 %%
%% Please do not use \input{...} to include other tex files.       %%
%% Submit your LaTeX manuscript as one .tex document.              %%
%%                                                                 %%
%% All additional figures and files should be attached             %%
%% separately and not embedded in the \TeX\ document itself.       %%
%%                                                                 %%
%%%%%%%%%%%%%%%%%%%%%%%%%%%%%%%%%%%%%%%%%%%%%%%%%%%%%%%%%%%%%%%%%%%%%

%%\documentclass[referee,sn-basic]{sn-jnl}% referee option is meant for double line spacing

%%=======================================================%%
%% to print line numbers in the margin use lineno option %%
%%=======================================================%%

%%\documentclass[lineno,pdflatex,sn-basic]{sn-jnl}% Basic Springer Nature Reference Style/Chemistry Reference Style

%%=========================================================================================%%
%% the documentclass is set to pdflatex as default. You can delete it if not appropriate.  %%
%%=========================================================================================%%

%%\documentclass[sn-basic]{sn-jnl}% Basic Springer Nature Reference Style/Chemistry Reference Style

%%Note: the following reference styles support Namedate and Numbered referencing. By default the style follows the most common style. To switch between the options you can add or remove �Numbered� in the optional parenthesis. 
%%The option is available for: sn-basic.bst, sn-chicago.bst%  
 
%%\documentclass[pdflatex,sn-nature]{sn-jnl}% Style for submissions to Nature Portfolio journals
%%\documentclass[pdflatex,sn-basic]{sn-jnl}% Basic Springer Nature Reference Style/Chemistry Reference Style
\documentclass[pdflatex,sn-mathphys-num]{sn-jnl}% Math and Physical Sciences Numbered Reference Style
%%\documentclass[pdflatex,sn-mathphys-ay]{sn-jnl}% Math and Physical Sciences Author Year Reference Style
%%\documentclass[pdflatex,sn-aps]{sn-jnl}% American Physical Society (APS) Reference Style
%%\documentclass[pdflatex,sn-vancouver-num]{sn-jnl}% Vancouver Numbered Reference Style
%%\documentclass[pdflatex,sn-vancouver-ay]{sn-jnl}% Vancouver Author Year Reference Style
%%\documentclass[pdflatex,sn-apa]{sn-jnl}% APA Reference Style
%%\documentclass[pdflatex,sn-chicago]{sn-jnl}% Chicago-based Humanities Reference Style

%%%% Standard Packages
%%<additional latex packages if required can be included here>

\usepackage{graphicx}%
\usepackage{multirow}%
\usepackage{amsmath,amssymb,amsfonts}%
\usepackage{amsthm}%
\usepackage{mathrsfs}%
\usepackage[title]{appendix}%
\usepackage{xcolor}%
\usepackage{textcomp}%
\usepackage{manyfoot}%
\usepackage{booktabs}%
\usepackage{algorithm}%
\usepackage{algorithmicx}%
\usepackage{algpseudocode}%
\usepackage{listings}%
%%%%

%%%%%=============================================================================%%%%
%%%%  Remarks: This template is provided to aid authors with the preparation
%%%%  of original research articles intended for submission to journals published 
%%%%  by Springer Nature. The guidance has been prepared in partnership with 
%%%%  production teams to conform to Springer Nature technical requirements. 
%%%%  Editorial and presentation requirements differ among journal portfolios and 
%%%%  research disciplines. You may find sections in this template are irrelevant 
%%%%  to your work and are empowered to omit any such section if allowed by the 
%%%%  journal you intend to submit to. The submission guidelines and policies 
%%%%  of the journal take precedence. A detailed User Manual is available in the 
%%%%  template package for technical guidance.
%%%%%=============================================================================%%%%

%% as per the requirement new theorem styles can be included as shown below
\theoremstyle{thmstyleone}%
\newtheorem{theorem}{Theorem}%  meant for continuous numbers
%%\newtheorem{theorem}{Theorem}[section]% meant for sectionwise numbers
%% optional argument [theorem] produces theorem numbering sequence instead of independent numbers for Proposition
\newtheorem{proposition}[theorem]{Proposition}% 
%%\newtheorem{proposition}{Proposition}% to get separate numbers for theorem and proposition etc.

\theoremstyle{thmstyletwo}%
\newtheorem{example}{Example}%
\newtheorem{remark}{Remark}%

\theoremstyle{thmstylethree}%
\newtheorem{definition}{Definition}%

\raggedbottom
%%\unnumbered% uncomment this for unnumbered level heads

\begin{document}

\title[Article Title]{Análisis Matemático, Complejidad Asintótica y Graficación}

%%=============================================================%%
%% GivenName	-> \fnm{Joergen W.}
%% Particle	-> \spfx{van der} -> surname prefix
%% FamilyName	-> \sur{Ploeg}
%% Suffix	-> \sfx{IV}
%% \author*[1,2]{\fnm{Joergen W.} \spfx{van der} \sur{Ploeg} 
%%  \sfx{IV}}\email{iauthor@gmail.com}
%%=============================================================%%

\author*[1,2]{\fnm{Navarrete} \sur{Alejandro}}\email{dennis.navarrete@epn.edu.ec}

\author[2,3]{\fnm{Crespo} \sur{Samuel}}\email{samuel.crespo@epn.edu.ec}
\equalcont{These authors contributed equally to this work.}

\author[1,2]{\fnm{Lopez} \sur{Esteban}}\email{jose.lopez05@epn.edu.ec}
\equalcont{These authors contributed equally to this work.}

\author[2,3]{\fnm{Duran} \sur{John}}\email{john.duran@epn.edu.ec}
\equalcont{These authors contributed equally to this work.}

\affil*[1]{\orgdiv{Ingenieria en Ciencias de Datos e Inteligencia Artificial}, \orgname{Escuela Politecnica Nacional}, \orgaddress{\street{Av. 12 de Octubre}, \city{Quito}, \postcode{171048}, \state{Pichincha}, \country{Ecuador}}}

%%==================================%%
%% Sample for unstructured abstract %%
%%==================================%%

\abstract{The abstract serves both as a general introduction to the topic and as a brief, non-technical summary of the main results and their implications. Authors are advised to check the author instructions for the journal they are submitting to for word limits and if structural elements like subheadings, citations, or equations are permitted.}

%%================================%%
%% Sample for structured abstract %%
%%================================%%

% \abstract{\textbf{Purpose:} The abstract serves both as a general introduction to the topic and as a brief, non-technical summary of the main results and their implications. The abstract must not include subheadings (unless expressly permitted in the journal's Instructions to Authors), equations or citations. As a guide the abstract should not exceed 200 words. Most journals do not set a hard limit however authors are advised to check the author instructions for the journal they are submitting to.
% 
% \textbf{Methods:} The abstract serves both as a general introduction to the topic and as a brief, non-technical summary of the main results and their implications. The abstract must not include subheadings (unless expressly permitted in the journal's Instructions to Authors), equations or citations. As a guide the abstract should not exceed 200 words. Most journals do not set a hard limit however authors are advised to check the author instructions for the journal they are submitting to.
% 
% \textbf{Results:} The abstract serves both as a general introduction to the topic and as a brief, non-technical summary of the main results and their implications. The abstract must not include subheadings (unless expressly permitted in the journal's Instructions to Authors), equations or citations. As a guide the abstract should not exceed 200 words. Most journals do not set a hard limit however authors are advised to check the author instructions for the journal they are submitting to.
% 
% \textbf{Conclusion:} The abstract serves both as a general introduction to the topic and as a brief, non-technical summary of the main results and their implications. The abstract must not include subheadings (unless expressly permitted in the journal's Instructions to Authors), equations or citations. As a guide the abstract should not exceed 200 words. Most journals do not set a hard limit however authors are advised to check the author instructions for the journal they are submitting to.}

\keywords{keywordadad, Keyword2, Keyword3, Keyword4}

%%\pacs[JEL Classification]{D8, H51}

%%\pacs[MSC Classification]{35A01, 65L10, 65L12, 65L20, 65L70}

\maketitle

\section{Introduction}\label{sec1}
En el desarrollo de software, la creación de soluciones efectivas va más allá de lograr que un programa arroje el resultado correcto. A medida que las aplicaciones escalan y el volumen de datos procesados aumenta, la eficiencia de los algoritmos se convierte en un factor crucial. Un código que funciona bien con pequeñas cantidades de información puede volverse inmanejable si su estructura no está optimizada. Por ello, el análisis de la complejidad algorítmica resulta fundamental, ya que nos permite predecir el comportamiento de un programa y evaluar su consumo de tiempo y recursos antes de implementarlo a gran escala.

En el marco de este Taller de Aprendizaje Práctico-Experimental, el presente documento tiene como objetivo aplicar los conceptos de eficiencia algorítmica para evaluar el rendimiento de tres fragmentos de código iterativos escritos en el lenguaje de programación Java. A través de este ejercicio, se busca establecer una relación clara entre el diseño interno de cada algoritmo y su respuesta ante el crecimiento continuo de los datos de entrada, variable que se representará matemáticamente como $n$.

Para llevar a cabo este estudio, el análisis de cada bloque de código se divide en tres etapas consecutivas. En la primera fase, se realiza una inspección detallada línea por línea para extraer la función de coste. Esto implica contabilizar las operaciones elementales que ejecuta la computadora, con el fin de construir paso a paso el modelo matemático $f(n)$ exacto. Una vez definida esta ecuación, la segunda fase consiste en la determinación de su límite asintótico superior. En este punto, se simplifica la función original descartando las constantes y los términos de menor impacto para poder clasificar la complejidad del algoritmo utilizando la notación Big O ($\mathcal{O}$).

Finalmente, el análisis analítico se complementa con una demostración empírica y visual mediante la graficación de los resultados obtenidos. Las curvas asintóticas resultantes de los tres ejercicios se proyectarán de manera simultánea en un único plano cartesiano, relacionando el tiempo de ejecución en el eje vertical frente al volumen de datos en el eje horizontal. Esta comparativa final servirá para evidenciar de forma práctica las distintas tasas de crecimiento de cada enfoque, permitiendo concluir visualmente cuál de los tres métodos ofrece el mejor desempeño a largo plazo.


\section{Results}\label{sec2}

El análisis de complejidad algorítmica realizado sobre las fracciones de código evaluados nos permitió determinar el comportamiento del tiempo de ejecución y el costo computacional en función del tamaño de la entrada.

Para la función A, el algoritmo implementa un bucle simple donde la variable utilizada para controlar este mismo incrementa de forma unitaria en cada iteración, lo cual da como resultado un número de operaciones directamente proporcional a n (tamaño de entrada), por lo que su complejidad es de $O(n)$.

Por otro lado, la función B introduce dos bucles anidados que recorren el rango completo de la entrada, en consecuencia, el número total de operaciones elementales se multiplicarán (n x n). Esto genera una función de costo cuadrática, estableciendo una complejidad algorítmica de $O(n^2)$.

Para finalizar, la función C implementa un bucle iterativo que se divide de manera sucesiva (actual = actual / 2), donde el valor de la variable utilizada para el bucle se va reduciendo a la mitad cada vez que itera, este comportamiento restringe al número total de ejecuciones a el número de veces que se puede dividir el tamaño de la entrada para 2, determinando una complejidad algorítmica de $O(\log n)$.

\section{Equations}\label{sec3}

\begin{lstlisting}[language=Java,caption={C\'odigo A}]
public void funcionA(int n) {
    int total = 0;
    for (int i = 0; i < n; i++) {
        total += i;
    }
    System.out.println(total);
}
\end{lstlisting}


\begin{enumerate}
    \item Conteo de operaciones elementales. 

Las operaciones elementales de este ejercicio son las siguientes:
\begin{itemize}
    \item \texttt{int total = 0}: $C_1 \times 1$
    \item \texttt{(int i = 0; i < n; i++)}: $C_2 \times n$
    \item \texttt{total += i}: $C_3 \times 1$
    \item \texttt{System.out.println(total)}: $C_4 \times 1$
\end{itemize}

    \item C\'alculo de la complejidad asint\'otica.
\end{enumerate}

\textit{Funci\'on A:}
\[
F(n) = C1 + (C2 * n) + C3 + C4
\]
\textit{Funci\'on A:}
\[
F(n) = C2 * n
\]
\textit{Funci\'on A:}
\[
F(n) = n
\]
\[
O(n)
\]

\begin{lstlisting}[language=Java,caption={C\'odigo B}]
public void funcionB(int n) {
    for (int i = 0; i < n; i++) {
        for (int j = 0; j < n; j++) {
            System.out.println("Procesando coordenadas:" + i + "," + j);
        }
    }
}
\end{lstlisting}


\begin{enumerate}
    \item Conteo de operaciones elementales. 

Las operaciones elementales de este ejercicio son las siguientes:
\begin{itemize}
    \item \texttt{(int i = 0; i < n; i++)}: $C_1 \times n$
    \item \texttt{(int j = 0; j < n; j++)}: $C_2 \times n$
    \item \texttt{System.out.println("Procesando coordenadas:" + i + "," + j)}: $C_3 \times 1$
\end{itemize}

    \item C\'alculo de la complejidad asint\'otica.
\end{enumerate}

\textit{Funci\'on B:}
\[
F(n) = (C1 * n) * (C2 * n) * C3 
\]
\textit{Funci\'on B:}
\[
F(n) = C1 * C2 * C3 * n^2
\]
\textit{Funci\'on B:}
\[
F(n) = n^2
\]
\[
O(n^2)
\]

\begin{lstlisting}[language=Java,caption={C\'odigo C}]
public void funcionC(int n) {
    int actual = n;
    while (actual > 1) {
        System.out.println("Valor:" + actual);
        actual = actual / 2;
    }
}
\end{lstlisting}
\begin{enumerate}
    \item Conteo de operaciones elementales. 

En este ejercicio, las operaciones elementales son las siguientes:
\begin{itemize}
    \item \texttt{int actual = n}: $C_1 \times 1$
    \item \texttt{while (actual > 1)}: $n$
    \item \texttt{System.out.println("Valor:" + actual)}: $C_2 \times 1$
    \item \texttt{actual = actual / 2}: $C_3 \rightarrow \log n$
\end{itemize}

    \item C\'alculo de la complejidad asint\'otica.
\end{enumerate}

\textit{Funci\'on C:}
\[
T(n) = C1 + n(C2 + \log n) = C1 + \log n = O(n)
\]
\textit{Funci\'on C:}
\[
T(n) = C1 + \log n
\]
\textit{Funci\'on C:}
\[
T(n) = \log n
\]
\[
O(\log n)
\]

\section{Graphs}\label{sec4}

\begin{figure}[h]
    \centering
    \includegraphics[width=0.6\textwidth]{grafico1.jpeg}
    \caption{Complejidades sobrepuestas de las funciones A, B y C.}
    \label{fig:Comparación de complejidades}
\end{figure}

La representación gráfica de las tres funciones confirma de manera visual los resultados teóricos obtenidos mediante el conteo de operaciones elementales. En el plano cartesiano, donde el eje horizontal representa el tamaño de entrada $n$ y el eje vertical representa el costo computacional, se puede observar claramente la diferencia en las tasas de crecimiento de cada algoritmo.

Para la funcion C, se nota un crecimiento inicial que con medida que avanza se aplana de forma casi horizontal, lo cual demuestra alta eficiencia en valores altos de entrada.

Por otro lado, la funcion A mantiene un crecimiento uniforme, tanto para valores altos como bajos, lo cual demuestra un crecimiento proporcional al número de entrada con el tiempo.

Pero, para la funcion B se mira un crecimiento acelerado de manera vertical y acelerada cuando el tamaño de entrada empieza a crecer, el comportamiento demuestra de manera visual el alto costo algorítmico. 

\section{Methods}\label{sec11}

Topical subheadings are allowed. Authors must ensure that their Methods section includes adequate experimental and characterization data necessary for others in the field to reproduce their work. Authors are encouraged to include RIIDs where appropriate. 

\textbf{Ethical approval declarations} (only required where applicable) Any article reporting experiment/s carried out on (i)~live vertebrate (or higher invertebrates), (ii)~humans or (iii)~human samples must include an unambiguous statement within the methods section that meets the following requirements: 

\begin{enumerate}[1.]
\item Approval: a statement which confirms that all experimental protocols were approved by a named institutional and/or licensing committee. Please identify the approving body in the methods section

\item Accordance: a statement explicitly saying that the methods were carried out in accordance with the relevant guidelines and regulations

\item Informed consent (for experiments involving humans or human tissue samples): include a statement confirming that informed consent was obtained from all participants and/or their legal guardian/s
\end{enumerate}

If your manuscript includes potentially identifying patient/participant information, or if it describes human transplantation research, or if it reports results of a clinical trial then  additional information will be required. Please visit (\url{https://www.nature.com/nature-research/editorial-policies}) for Nature Portfolio journals, (\url{https://www.springer.com/gp/authors-editors/journal-author/journal-author-helpdesk/publishing-ethics/14214}) for Springer Nature journals, or (\url{https://www.biomedcentral.com/getpublished/editorial-policies\#ethics+and+consent}) for BMC.

\section{Discussion}\label{sec12}

El estudio de las tres funciones de código A, B y C demuestra como la estructura de un algoritmo determina su crecimiento asintotico y su coste computacional. Aunque todas reciben como parametro de entrada \(n\), sus costes de ejecución cambian a diversos ritmos en base al crecimiento del tamaño de entrada. Tal comparación deja ver lo importante que es la complejidad algoritmica al evaluar la escalabilidad del software.

El ejercicio A muestra una complejidad lineal temporal \(O(n)\), dado que su operación principal es ejecutada una vez por iteración de bucle que recorre el rango de entrada. Su coste computacional se incrementa de manera proporcional al tamaño de la entrada.

El ejercicio B presenta una complejidad temporal cuadrática \(O(n^2)\), ya que cada uno de sus bucles anidados realiza \(n\) iteraciones. Entonces, el número de ejecuciones de ejecuciones de la instrucción más interna aumenta de forma cuadrática, logrando que esta función sea más costosa para entradas enormes.

El ejercicio C demuestra una complejidad temporal algorítmica \(O(\log n)\), porque el valor de su variable de control se divide para dos en cada iteración, tal reducción, repetidas veces, disminuye de manera sustancial la cantidad de iteraciones necesarias a medida que \(n\) aumenta.
Para tamaños de entrada lo suficientemente grandes, la función cuadrática crece más rápido en comparación a la lineal, mientras que la función logaritmica crece, de manera considerable, más despacio que ambas. Siendo la función C la que escala de manera más eficiente en base a la cantidad de iteraciones. La función B es la más sensible al aumento del tamaño de la entrada. Y la función A tiene un incremento predecible en la carga de trabajo computacional.
La notación \(Big O\) describe una cota superior asintótica en lugar del tiempo de ejecución de un programa. Los detalles de implementación, características del hardware, coste de las operaciones y los valores de entrada son aspectos de los cuales depende el tiempo de ejecución. El análisis asintótico da un marco teórico necesario para comparar algoritmos de forma independiente de un entorno de ejecución concreto.
\section{Conclusion}\label{sec13}

El desarrollo de este taller práctico permitió evidenciar de forma concreta la profunda influencia que tiene la estructura interna de un algoritmo sobre su rendimiento computacional. A través del cálculo detallado de las funciones de coste y su posterior simplificación asintótica, se demostró que el diseño lógico del código es tan crítico como su funcionalidad. Los tres fragmentos analizados, aunque sintácticamente breves en Java, exhibieron comportamientos radicalmente distintos frente al aumento del volumen de datos, reafirmando que no todas las soluciones lógicamente válidas 
son igualmente viables cuando se llevan a una escala mayor.
Al contrastar los resultados matemáticos, se identificaron tres categorías fundamentales de complejidad que marcan la diferencia en el diseño de software. La Función A, estructurada mediante un ciclo iterativo simple, presentó un crecimiento estrictamente lineal $\mathcal{O}(n)$, lo que indica que su tiempo de ejecución aumenta de manera proporcional a los datos de entrada. Por el contrario, la presencia de iteraciones anidadas en la Función B desencadenó un crecimiento cuadrático $\mathcal{O}(n^2)$, demostrando ser el enfoque más ineficiente de los tres y el más susceptible a 
generar cuellos de botella computacionales ante grandes volúmenes de información. En el extremo opuesto, la Función C destacó por su eficiencia al dividir sucesivamente el tamaño del problema, logrando un tiempo de ejecución logarítmico $\mathcal{O}(\log n)$ que garantiza una escalabilidad excepcional.
Finalmente, el análisis gráfico consolidó estas abstracciones teóricas en una comprobación empírica y visual. Al proyectar las tres curvas de rendimiento de manera simultánea en el plano cartesiano, quedó ilustrada con total claridad la separación de los límites asintóticos. Mientras que la curva logarítmica tendía a estabilizarse rápidamente y la curva lineal mantenía un ritmo predecible, la curva cuadrática experimentó un crecimiento acelerado e insostenible en el eje del tiempo. En definitiva, este ejercicio subraya la importancia de la notación Big O como una herramienta analítica indispensable, la cual permite a los desarrolladores prever el impacto de sus decisiones de código y optimizar el consumo de los recursos del sistema antes de su implementación final.

\section*{Declarations}

Some journals require declarations to be submitted in a standardised format. Please check the Instructions for Authors of the journal to which you are submitting to see if you need to complete this section. If yes, your manuscript must contain the following sections under the heading `Declarations':

\begin{itemize}
\item Funding
\item Conflict of interest/Competing interests (check journal-specific guidelines for which heading to use)
\item Ethics approval and consent to participate
\item Consent for publication
\item Data availability 
\item Materials availability
\item Code availability 
\item Author contribution
\end{itemize}

\noindent
Adem\'as se declara el uso de plantilla de Springer Nature para la redacción del presente documento, la cual se encuentra disponible en :

%%===================================================%%
%% For presentation purpose, we have included        %%
%% \bigskip command. Please ignore this.             %%
%%===================================================%%
\bigskip
\begin{flushleft}%
Editorial Policies for:

\bigskip\noindent
Springer journals and proceedings: \url{https://www.springer.com/gp/editorial-policies}

\bigskip\noindent
Nature Portfolio journals: \url{https://www.nature.com/nature-research/editorial-policies}
\end{flushleft}


\begin{thebibliography}{99}

\bibitem{bhargava2016}
Bhargava, A. Y. (2016). \textit{Grokking Algorithms: An illustrated guide for programmers and other curious people}. Manning Publications.

\bibitem{lafore2017} 
Lafore, R. (2017). \textit{Data Structures and Algorithms in Java} (2\textsuperscript{da} ed.). Sams Publishing.

\bibitem{astudillo2026}
Galarza, D. (2026). \textit{Análisis Matemático y Complejidad Asintótica} [Diapositivas de clase]. Asignatura de Estructura de Datos y Algoritmos, Escuela Politécnica Nacional.

\end{thebibliography}

\end{document}
